\documentclass[11pt,a4paper]{article}
\usepackage[italian]{babel}
\usepackage[margin=2.5cm]{geometry}
\usepackage{parskip}
\usepackage{amsmath, amssymb, amsthm}
\usepackage{booktabs}
\usepackage{array}
\usepackage{xcolor}
\usepackage{needspace}
\usepackage{enumitem}
\usepackage{listings}
\usepackage{tikz}
\usepackage[most]{tcolorbox}
\usepackage[hidelinks]{hyperref}
\usetikzlibrary{decorations.pathreplacing, shapes.geometric, arrows.meta, positioning, automata}

% Niente vedove/orfane: i paragrafi non lasciano righe isolate a fine/inizio pagina
\widowpenalty=10000
\clubpenalty=10000

% Elenchi più compatti
\setlist{itemsep=2pt, parsep=0pt, topsep=4pt, partopsep=0pt}

% --- Riquadri con bordo nero, senza numero (si spezzano tra due pagine) ---
% Uso: \begin{definizione} ... \end{definizione}
%      \begin{teorema}[Nome] ... \end{teorema}  (il nome è facoltativo)
\tcbset{nota/.style={enhanced jigsaw, breakable, colback=white,
  colframe=black, boxrule=0.5pt, sharp corners}}
\NewTColorBox{definizione}{}{nota,
  before upper={\textbf{Definizione.}\ }}
\NewTColorBox{teorema}{o}{nota,
  before upper={\textbf{Teorema\IfValueT{#1}{ (#1)}.}\ }}
\NewTColorBox{proposizione}{o}{nota,
  before upper={\textbf{Proposizione\IfValueT{#1}{ (#1)}.}\ }}
\NewTColorBox{proprieta}{o}{nota,
  before upper={\textbf{Proprietà\IfValueT{#1}{ (#1)}.}\ }}
\NewTColorBox{assioma}{o}{nota, boxrule=1pt,
  before upper={\textbf{Assioma\IfValueT{#1}{ (#1)}.}\ }}

% --- Ambienti semplici (senza riquadro) ---
% Il titolo sta in cima al blocco, anche se il contenuto inizia con un riquadro o
% con codice; e non resta mai da solo in fondo a una pagina.
% Uso: \begin{esempio}[Titolo facoltativo] ... \end{esempio}
\newcommand{\newrunin}[2]{%
  \NewDocumentEnvironment{#1}{o}{%
    \par\Needspace{4\baselineskip}\addvspace{\medskipamount}\noindent
    \textbf{#2}\IfValueT{##1}{ (##1)}\textbf{.}\ \ignorespaces}%
  {\par\addvspace{\medskipamount}}}
\newrunin{esempio}{Esempio}
\newrunin{esercizio}{Esercizio}
\NewTColorBox{osservazione}{}{nota,
  before upper={\textbf{Osservazione.}\ }}

% V e F nelle tavole di verità
\newcommand{\V}{\textsf{V}}
\newcommand{\F}{\textsf{F}}

% --- Codice C ---
% Uso: \begin{codice} ... \end{codice}      codice C numerato
%      \begin{comando} ... \end{comando}    comandi da terminale
%      \begin{risultato} ... \end{risultato} output di un programma
%      \lstinline|printf("ciao");|           codice dentro al testo
\lstdefinestyle{c}{
  language=C,
  basicstyle=\ttfamily\small,
  keywordstyle=\color{blue}\bfseries,
  commentstyle=\color{gray}\itshape,
  stringstyle=\color{red!70!black},
  numbers=left,
  numberstyle=\tiny\color{gray},
  numbersep=8pt,
  columns=flexible,
  keepspaces=true,
  breaklines=true,
  showstringspaces=false,
  tabsize=4,
  morekeywords={include, define},
  % lettere accentate nei commenti
  literate={à}{{\`a}}1 {è}{{\`e}}1 {é}{{\'e}}1 {ì}{{\`i}}1
           {ò}{{\`o}}1 {ù}{{\`u}}1 {È}{{\`E}}1 {À}{{\`A}}1
}
\lstset{style=c}
\newtcblisting{codice}{listing only, nota, left=6mm,
  listing options={style=c}}
\newtcblisting{comando}{listing only, nota,
  listing options={basicstyle=\ttfamily\small, numbers=none, language={},
    columns=flexible, keepspaces=true, breaklines=true}}
\newtcblisting{risultato}{listing only, enhanced jigsaw, breakable,
  colback=black!5, colframe=black, boxrule=0.5pt, sharp corners,
  listing options={basicstyle=\ttfamily\small, numbers=none, language={},
    columns=flexible, keepspaces=true, breaklines=true}}

% --- Diagrammi di flusso (simboli standard) ---
\tikzset{
  terminale/.style = {ellipse, draw, minimum width=2.5cm, minimum height=0.9cm},
  io/.style        = {trapezium, trapezium left angle=70, trapezium right angle=110,
                      draw, minimum width=2.5cm, minimum height=0.9cm},
  processo/.style  = {rectangle, draw, minimum width=2.5cm, minimum height=0.9cm},
  decisione/.style = {diamond, draw, aspect=2, minimum width=2.5cm},
  freccia/.style   = {thick, -{Stealth}}
}

\title{Mappe di Karnaugh: il punto di vista del libro}
\author{Marco Cerbella}
\date{Lezione 6}

\begin{document}
\maketitle

\section{Mappe di Karnaugh: il punto di vista del libro}

Questa lezione riprende le mappe di Karnaugh (lezione 5) seguendo il libro di testo di Harris \& Harris, \emph{Digital Design and Computer Architecture}, sezione 2.7.

Dopo aver provato a minimizzare equazioni booleane con l'algebra di Boole, ci si accorge che, se non si sta attenti, si può finire con un'equazione completamente diversa invece che semplificata. Le mappe di Karnaugh (\emph{K-map}) sono un metodo grafico per semplificare le equazioni booleane; furono inventate nel 1953 da Maurice Karnaugh, ingegnere delle telecomunicazioni ai Bell Labs. Funzionano bene per problemi con fino a quattro variabili e, soprattutto, aiutano a capire come manipolare le equazioni booleane.

\begin{osservazione}
Maurice Karnaugh (1924--) si è laureato in fisica al City College di New York nel 1948 e ha conseguito il dottorato in fisica a Yale nel 1952. Ha lavorato ai Bell Labs e all'IBM dal 1952 al 1993 e come professore di informatica al Polytechnic University di New York dal 1980 al 1999.
\end{osservazione}

\section{Perché una mappa}

Minimizzare una funzione logica significa combinare dei termini. Due termini che contengono un implicante $P$ e la forma vera e quella complementare di una stessa variabile $A$ si combinano eliminando $A$:
\[
PA + PA' = P.
\]
Le mappe di Karnaugh rendono evidenti questi termini combinabili, mettendoli uno accanto all'altro in una griglia.

\subsection{Mappa di una funzione a tre ingressi}

Consideriamo la funzione $Y$ di tre ingressi $A, B, C$ con $Y = 1$ solo per $ABC = 000$ e $ABC = 001$ (e $0$ altrove). La riga superiore della mappa riporta i quattro possibili valori degli ingressi $A$ e $B$; la colonna di sinistra riporta i due possibili valori dell'ingresso $C$. Ogni quadrato della mappa corrisponde a una riga della tabella di verità e contiene il valore dell'uscita $Y$ per quella riga.

\begin{center}
\begin{tabular}{c|cccc}
 & \multicolumn{4}{c}{$AB$}\\
$C$ & $00$ & $01$ & $11$ & $10$ \\
\midrule
$0$ & 1 & 0 & 0 & 0 \\
$1$ & 1 & 0 & 0 & 0 \\
\end{tabular}
\end{center}

Come ogni riga di una tabella di verità, ogni quadrato di una K-map rappresenta un singolo \emph{mintermine}: la cella in alto a sinistra è $A'B'C'$, quella sotto è $A'B'C$, e così via per tutti gli otto mintermini.

\section{Codice di Gray e adiacenza}

Ogni quadrato (mintermine) differisce da un quadrato adiacente per una sola variabile. Quindi i quadrati adiacenti condividono tutti i letterali tranne uno, che compare in forma vera in un quadrato e in forma complementare nell'altro. Per esempio i quadrati che rappresentano i mintermini $A'BC'$ e $A'BC$ sono adiacenti e differiscono solo per la variabile $C$.

Le combinazioni di $A$ e $B$ nella riga superiore sono in un ordine particolare: $00, 01, 11, 10$. Questo ordine si chiama \emph{codice di Gray}. Si distingue dal normale ordine binario ($00, 01, 10, 11$) perché voci adiacenti differiscono per una sola variabile: per esempio passando da $01$ a $11$ cambia solo $A$ (da $0$ a $1$), mentre passando da $01$ a $10$ cambierebbe sia $A$ (da $1$ a $0$) sia $B$ (da $0$ a $1$). Scrivere le combinazioni in ordine binario non darebbe la proprietà desiderata, cioè quadrati adiacenti che differiscono per una sola variabile.

\begin{osservazione}
I codici di Gray valgono per un numero qualunque di bit. Per esempio una sequenza di Gray a $3$ bit è
\[
000,\ 001,\ 011,\ 010,\ 110,\ 111,\ 101,\ 100.
\]
Furono brevettati nel 1953 (brevetto USA 2\,632\,058) da Frank Gray, ricercatore ai Bell Labs; sono utili soprattutto negli encoder meccanici, perché un leggero disallineamento provoca un errore in un solo bit. Un problema collegato fu proposto da Lewis Carroll nel 1879: collegare due parole della stessa lunghezza con parole intermedie che differiscono dalla precedente per una sola lettera, per esempio da SHIP a DOCK: SHIP, SLIP, SLOP, SLOT, SOOT, LOOT, LOOK, LOCK, DOCK.
\end{osservazione}

\subsection{La mappa si ``arrotola''}

La K-map si ``avvolge'' su se stessa: i quadrati all'estrema destra sono di fatto adiacenti a quelli all'estrema sinistra, perché differiscono per una sola variabile, $A$. In altre parole si può arrotolare la mappa in un cilindro, poi unire le estremità del cilindro per ottenere un toro (una ciambella) e avere comunque la garanzia che quadrati adiacenti differiscano per una sola variabile.

\section{Minimizzazione con le K-map}

\subsection{Ragionare per cerchi}

Nella K-map della funzione precedente nell'equazione compaiono solo due mintermini, $A'B'C'$ e $A'B'C$, indicati dagli $1$ nella colonna di sinistra. Leggere i mintermini da una K-map è esattamente equivalente a leggere le equazioni in somma di prodotti direttamente dalla tabella di verità. Come prima si può usare l'algebra di Boole:
\[
Y = A'B'C' + A'B'C = A'B'(C' + C) = A'B'.
\]
Le K-map permettono di fare questa semplificazione graficamente, cerchiando gli $1$ in quadrati adiacenti. Per ogni cerchio si scrive l'\emph{implicante} corrispondente. Un implicante è il prodotto di uno o più letterali; le variabili le cui forme vera e complementare sono entrambe dentro il cerchio vengono escluse dall'implicante. Nel nostro caso la variabile $C$ ha dentro il cerchio sia la forma vera ($1$) sia la forma complementare ($0$), quindi non viene inclusa nell'implicante: $Y$ è vera quando $A = B = 0$, indipendentemente da $C$. L'implicante è $A'B'$, lo stesso risultato ottenuto con l'algebra di Boole.

\subsection{Regole}

La K-map dà un modo visivo semplice per minimizzare la logica: basta cerchiare tutti i blocchi rettangolari di $1$ della mappa, usando il minor numero possibile di cerchi. Ogni cerchio deve essere il più grande possibile. Poi si leggono gli implicanti cerchiati.

Più formalmente, un'equazione booleana è minimizzata quando è scritta come somma del minor numero di \emph{implicanti primi}. Ogni cerchio sulla K-map rappresenta un implicante; i cerchi più grandi possibili sono gli implicanti primi.

\begin{definizione}
Un \emph{implicante primo} è un implicante che non può essere ulteriormente ingrandito, cioè non può essere combinato con altri mintermini per eliminare una variabile.
\end{definizione}

Per esempio, nella mappa precedente $A'B'C'$ e $A'B'C$ sono implicanti, ma non primi. Solo $A'B'$ è un implicante primo.

Le regole per trovare un'equazione minimizzata da una K-map sono:
\begin{itemize}
    \item usare il minor numero di cerchi necessario a coprire tutti gli $1$;
    \item tutti i quadrati dentro ogni cerchio devono contenere $1$;
    \item ogni cerchio deve coprire un blocco rettangolare che ha, in ogni direzione, un numero di quadrati che è una potenza di $2$ ($1$, $2$ o $4$);
    \item ogni cerchio deve essere il più grande possibile;
    \item un cerchio può passare oltre i bordi della K-map (avvolgendosi);
    \item un $1$ in una K-map può essere cerchiato più volte, se questo consente di usare meno cerchi.
\end{itemize}

\begin{esempio}[Minimizzazione di una funzione di tre variabili]
Sia $Y = F(A,B,C)$ con la K-map seguente:
\begin{center}
\begin{tabular}{c|cccc}
 & \multicolumn{4}{c}{$AB$}\\
$C$ & $00$ & $01$ & $11$ & $10$ \\
\midrule
$0$ & 1 & 0 & 1 & 1 \\
$1$ & 1 & 0 & 0 & 1 \\
\end{tabular}
\end{center}
Si cerchiano gli $1$ con il minor numero di cerchi possibile. Ogni cerchio rappresenta un implicante primo, e le dimensioni di ciascun cerchio sono potenze di due ($2 \times 1$ e $2 \times 2$). Si forma l'implicante primo di ciascun cerchio scrivendo le variabili che compaiono nel cerchio solo in forma vera o solo in forma complementare.

\begin{itemize}
    \item Il cerchio $2 \times 1$ comprende i quadrati $AB = 11$ e $AB = 10$ nella riga $C = 0$. In esso compaiono sia $B$ sia $B'$, quindi $B$ non fa parte dell'implicante; compaiono invece solo $A$ in forma vera e $C$ in forma complementare: l'implicante primo è $AC'$.
    \item Il cerchio $2 \times 2$ copre tutti i quadrati in cui $B = 0$ (le colonne $00$ e $10$, che sono adiacenti perché la mappa si avvolge ai bordi): l'implicante primo è $B'$.
\end{itemize}
Perciò
\[
Y = AC' + B'.
\]
Si noti che il quadrato in alto a destra ($AB = 10$, $C = 0$) è coperto due volte, per rendere gli implicanti primi i più grandi possibili: come nelle tecniche con l'algebra di Boole, equivale a condividere un mintermine per ridurre la dimensione di un implicante. Si noti anche che il cerchio da quattro quadrati si avvolge attorno ai lati della K-map.
\end{esempio}

\begin{osservazione}
Il libro prosegue con l'esempio~2.10, un decodificatore per display a sette segmenti: un circuito con ingresso a $4$ bit $D_{3:0}$ e sette uscite $S_a, \dots, S_g$ che pilotano i LED per mostrare una cifra da $0$ a $9$ (gli ingressi illegali $10$--$15$ producono un display spento). Se ne cerca la tabella di verità e, con le K-map, le equazioni delle uscite $S_a$ e $S_b$. Le pagine fornite con questa lezione terminano all'inizio dell'esempio: vi si ricorda in particolare di etichettare righe e colonne in ordine di Gray anche quando si riportano i valori dell'uscita, e che l'insieme minimo di implicanti primi non è unico.
\end{osservazione}

\section{Riferimenti}

Harris \& Harris, \emph{Digital Design and Computer Architecture} (in italiano \emph{Sistemi digitali e architettura dei calcolatori}), sezione 2.7, disponibile su UniStudium (pagine 71--75 fornite). Per le slide della lezione si veda la lezione~5.

\end{document}
